% Original analytic teaching diagram, 2026-11-12.
% Implications between convergence modes on one probability space (1 <= p < r < infinity).
% Solid arrows: always valid. Dashed arrows: valid under the stated extra condition.
\definecolor{VB}{HTML}{7188A3}\definecolor{VBF}{HTML}{E5EAF0}
\definecolor{VG}{HTML}{829C8B}\definecolor{VGF}{HTML}{E8EEE9}
\definecolor{VY}{HTML}{B4A16F}\definecolor{VYF}{HTML}{F1ECDD}
\definecolor{VR}{HTML}{A7787D}\definecolor{VRF}{HTML}{F0E5E6}
\definecolor{VP}{HTML}{9588A5}\definecolor{VPF}{HTML}{ECE8F0}
\definecolor{VI}{HTML}{3B4252}\definecolor{VM}{HTML}{6B7280}\definecolor{VL}{HTML}{4C4D4F}
\begin{tikzpicture}[x=1mm,y=1mm,>=Stealth,font=\figurefont\fontsize{8}{10}\selectfont,
lab/.style={text=VM,font=\figurefont\fontsize{7.5}{9}\selectfont},
mode/.style={draw=VB,fill=VBF,rounded corners=1.8pt,line width=.7pt,minimum width=22mm,minimum height=7.5mm,text=VI},
imp/.style={->,draw=VL,line width=.8pt},
cond/.style={->,draw=VP,line width=.8pt,dashed}]
\node[mode] (as) at (0,0) {几乎必然收敛};
\node[mode] (pr) at (42,0) {依概率收敛};
\node[mode] (lp) at (84,0) {$L^{p}$收敛};
\node[mode] (lr) at (84,-22) {$L^{r}$收敛};
\draw[imp] (as.east) -- (pr.west);
\draw[imp] (lp.west) -- (pr.east);
\draw[imp] (lr.north) -- (lp.south);
\node[lab,right] at (85,-11) {$r>p$};
\draw[cond] ([xshift=-4mm]pr.north) to[out=135,in=45] node[lab,above,pos=.5] {存在子列} ([xshift=4mm]as.north);
\draw[cond] ([xshift=4mm]pr.north) to[out=45,in=135] node[lab,above,pos=.5] {$\lvert X_n\rvert^{p}$一致可积} ([xshift=-4mm]lp.north);
\draw[cond] (as.south) to[out=-60,in=-120,looseness=.9] node[lab,below,pos=.5] {可积上界支配$\lvert X_n\rvert^{p}$} ([xshift=-7mm]lp.south);
\end{tikzpicture}
